\section{Bottleneck Identification}

\subsection{ATRI Top 50: Geographic Distribution}

Table~\ref{tab:atri-top10} shows the 10 highest-cost bottleneck locations from the 2024 ATRI report. The Atlanta metropolitan corridor dominates the top of the list, with four locations in the top 10: I-285/I-20 (\$916M), I-285/I-85 (\$789M), I-75/I-285 (\$687M), and I-75/I-85 downtown connector (\$617M). Combined, the Atlanta corridor cluster accounts for \$3.0 billion in annual freight congestion cost --- more than the entire New England bottleneck inventory.

\begin{table}[h!]
\centering
\caption{ATRI Top 10 Freight Bottlenecks by Annual Cost, 2024}
\label{tab:atri-top10}
\begin{tabular}{rlrrr}
\toprule
\textbf{Rank} & \textbf{Location} & \textbf{Route} & \textbf{Tier} & \textbf{Cost (\$M/yr)} \\
\midrule
1 & I-285/I-20 Atlanta & I-285 & T2 & 916 \\
2 & I-95/SR-4 Fort Lee NJ & I-95 & T1 & 848 \\
3 & I-95/I-695 Baltimore & I-95 & T1 & 812 \\
4 & I-285/I-85 Atlanta & I-285 & T2 & 789 \\
5 & I-95/I-476 Philadelphia & I-95 & T1 & 762 \\
6 & I-10/I-110 Los Angeles & I-10 & T1 & 731 \\
7 & I-270/I-70 St. Louis & I-270 & T2 & 698 \\
8 & I-75/I-285 Atlanta & I-75 & T1 & 687 \\
9 & I-95/I-395/I-695 Baltimore & I-95 & T1 & 654 \\
10 & I-80/I-94 Chicago & I-80 & T1 & 641 \\
\midrule
\multicolumn{4}{r}{\textbf{Top 10 total}} & \textbf{\$7,538M} \\
\bottomrule
\end{tabular}
\end{table}

\subsection{Corridor-Level A1 Scores vs. ATRI Rankings}

The ROUTE A1 (Throughput Gap) dimension measures 90th-percentile AADT as a proxy for congestion. We test whether A1 scores predict ATRI bottleneck density.

From the v1.1 corpus of 227 corridors, 14 corridors score A1 $\geq$ 6.0 (indicating chronic peak-period congestion). Of these 14 corridors, 11 (79\%) have at least one ATRI top-100 location. The correlation between corridor A1 score and ATRI bottleneck density is Spearman $\rho = 0.67$ ($p < 0.001$), confirming that A1 is a meaningful predictor of bottleneck location density.

However, the relationship is not linear. Corridors with A1 = 8--10 (severely congested) have similar bottleneck density to corridors with A1 = 6--7, because severe congestion tends to be localized to specific interchange clusters rather than distributed along the corridor. The A1 score captures corridor-level average stress; ATRI captures specific high-cost locations.

\subsection{ROUTE Score and ATRI Rank: Key Patterns}

Three patterns emerge from the ATRI-ROUTE integration:

\textbf{Pattern 1: Atlanta concentration.} Atlanta accounts for 4 of the top 10 and 7 of the top 50 ATRI locations. The I-285 beltway (T2) and its intersections with I-75, I-85, and I-20 form the most concentrated bottleneck cluster in the US. The I-285 corridor scores A1 = 8.2 in the ROUTE corpus --- correctly identified as highly congested --- but its T2 tier classification reflects its regional (not national) network role.

\textbf{Pattern 2: Northeast Corridor density.} I-95 has 9 of the top 50 ATRI locations, more than any other single corridor. As the only T1 corridor serving the Boston-Washington megalopolis, I-95 combines nationally-critical freight with the highest-density passenger vehicle traffic in the country. Its A1 score of 5.2 (moderate) understates its bottleneck density because the corridor's total mileage (1,919 miles) dilutes the extreme congestion concentrated in the New York-to-Baltimore segment (~150 miles).

\textbf{Pattern 3: I-40 absence.} I-40, despite scoring V/C = 0.84 (at target) in the ROUTE analysis, has zero ATRI top-50 locations. This confirms the finding from A.1: I-40 is the one T1 corridor where managed freight lanes are not urgently needed for bottleneck relief. Its investment priority is maintenance and EV charging, not capacity.
